\section{模型评价与推广}

\subsection{模型优点}

\begin{enumerate}[leftmargin=2.2em,itemsep=2pt]
  \item {\heiti 一体化框架：}问题一与问题二共用同一检测模型，判别结果与检测结果统计一致，避免了多模型之间的口径冲突，也节省了训练与部署成本；
  \item {\heiti 开放集判别能力：}基于极值理论的 EVT 判别模型在不构造大规模负样本的前提下完成二分类，数学依据充分、计算开销可忽略，且判别阈值具有明确的统计含义；
  \item {\heiti 针对性数据策略：}600 幅负样本训练显著抑制背景误检，并在问题一判别中取得最高的置信度分布分离度（$k=3.064$）；消融实验对增益来源作出严谨归因（提升主要来自骨干升级，Copy-Paste 因数据无分割掩膜无显著贡献）；
  \item {\heiti 工程可用性：}最终模型约 11.13 M 参数，端到端推理 8.87 ms/幅，满足港口集装箱实时巡检需求；全部实验固定随机种子，结果可复现。
\end{enumerate}

\subsection{模型缺点}

\begin{enumerate}[leftmargin=2.2em,itemsep=2pt]
  \item Rusty（锈蚀）类召回率偏低（0.282），大面积弱纹理缺陷与背景相似是当前主要性能瓶颈；
  \item 由于验证集不含无缺陷图像，问题一分类的混淆矩阵与 AUC 未能直接在本地量化，仅通过分布可分性间接验证；
  \item 赛后补充消融表明，Copy-Paste 增强在本数据集（无分割掩膜）上未产生可观测增益，原"改进"实验的增益主要来自骨干网络升级；此外 Rusty 类加权与 WD-Focal Loss 的端到端验证均为负结果，说明类别不平衡问题难以仅靠损失层改进解决；
  \item 鲁棒性扰动测试显示模型对高斯噪声敏感（$\sigma=50$ 时 mAP@0.5:0.95 由 0.204 降至 0.004），像素级噪声是当前鲁棒性的薄弱环节。
\end{enumerate}

\subsection{模型推广}

针对上述不足，可从以下方向进一步推广与改进：

\begin{enumerate}[leftmargin=2.2em,itemsep=2pt]
  \item {\heiti 损失函数改进迭代：}WD-Focal 端到端验证（负结果）表明其增益受损失尺度失衡制约，后续应在分类分支损失配平、$\gamma$ 与 $\alpha$ 限幅范围内重新调参，而非直接替换；
  \item {\heiti 小目标检测增强（已实测）：}P2 高分辨率检测头使 Hole 类 mAP@0.5:0.95 由 0.190 升至 0.210（+0.021）、召回率由 0.407 升至 0.453；多尺度测试增强（TTA）使整体 mAP@0.5:0.95 由 0.204 升至 0.215（+0.011），无需重新训练，建议优先采用；Soft-NMS 可进一步缓解密集目标抑制；
  \item {\heiti 缺陷分割扩展：}在检测基础上引入实例分割分支，输出破损轮廓与面积，为破损严重程度评估提供更细粒度信息；
  \item {\heiti 轻量化部署：}通过通道剪枝、知识蒸馏与量化将模型压缩至边缘设备可运行规模，结合多视角摄像头实现集装箱全表面自动巡检；
  \item {\heiti 评估完善：}扰动鲁棒性测试已补齐（噪声、亮度、模糊六类扰动）；后续在含真实负样本的标注集上重标定问题一阈值并计算混淆矩阵与 ROC-AUC，形成完整的模型检验报告。
\end{enumerate}

综上，本文构建的“YOLOv8 检测 + EVT 判别 + 多维度评估”方案，在数据存在类别不平衡、小目标密集、无负样本与复杂背景等现实困难的条件下，实现了准确、实时、可解释的集装箱破损检测，具备良好的工程推广价值。
